\chapter*{Plain-Language Summary}
\addcontentsline{toc}{chapter}{Plain-Language Summary}

\paragraph{The problem.}
Every large organisation --- a bank, a hospital, a university --- runs
computer networks that move millions of small data packets every minute. A
few of those packets belong to attackers. The job of a security team is to
spot the malicious traffic in the flood of normal traffic, which is
impossible by hand at scale. Security teams therefore deploy automated
intrusion detection systems (IDS) that read the traffic and raise alerts,
and increasingly those systems are built with machine learning trained on
historical labelled examples.

\paragraph{What was found.}
The most important finding is that the published numbers in many papers
are too good to be true. Those papers split their data randomly, which
lets the model see future traffic during training. When the same models
are tested with a more realistic protocol --- train on Monday to Wednesday,
test on Friday --- the macro F1 score drops from 0.86 to 0.44 on the
fifteen-way attack classification task. Macro F1 is a balanced version of
accuracy that gives each attack class equal weight, so it is the right
metric for a problem where benign traffic dominates. However, simplifying
the question to ``is this an attack or not?'' (without naming the specific
attack) recovers most of the lost performance: random forest still reaches
an F1 score of 0.86 on the realistic test.

\paragraph{What this means.}
A company that deploys an IDS based only on high benchmark scores from
random stratified splits may be disappointed in production. This thesis
shows the honest baseline, identifies which model (random forest) is most
reliable for deployment, and shows how to convert raw model output into
structured alerts tagged with MITRE ATT\&CK technique IDs. The thesis
also demonstrates that the model can be fooled by an attacker who makes
small adjustments to packet timing, and that a naïve attempt at fixing
this vulnerability actually makes things worse. The broader message:
machine-learning detectors that work in the lab do not always work in the
real world; the way the system is tested matters more than the model that
is chosen.
